\chapter{Somas de variáveis aleatórias e a lei dos grandes números}\label{ch:g12:sums}

Por que os cassinos sempre ganham no fim, e por que as pesquisas de
opinião funcionam? Porque as \omterm{def:g11:stat:mean}{médias} de muitas grandezas aleatórias
\omterm{def:g12:sums:indep}{independentes} flutuam cada vez menos. Este capítulo demonstra isso:
linearidade da \omterm{def:g12:randvar:exp}{esperança}, aditividade da \omterm{def:g12:randvar:exp}{variância} para variáveis
\omterm{def:g12:sums:indep}{independentes}, desigualdade de Bienaymé--Chebyshev e lei dos grandes
números.

\section{Somas de variáveis aleatórias}

Dadas duas variáveis aleatórias $X, Y$ sobre o mesmo \omterm{def:g11:prob:model}{espaço amostral}
finito $\Omega$, a soma $X + Y$ é a \omterm{def:g12:randvar:rv}{variável aleatória}
$\omega \mapsto X(\omega) + Y(\omega)$.

\begin{theorem}[Linearidade da esperança]\label{thm:g12:sums:linearity}
\index{esperança!linearidade}
Para todas as variáveis aleatórias $X, Y$ sobre $\Omega$ e todos
$a, b \in \R$:
\[
\E(X + Y) = \E(X) + \E(Y), \qquad \E(aX + b) = a\E(X) + b .
\]
\emph{Nenhuma hipótese de \omterm{def:g12:condprob:indep}{independência} é necessária.}
\end{theorem}

\begin{proof}
Escreva a \omterm{def:g12:randvar:exp}{esperança} como soma sobre os resultados: como
$\P(X = x) = \sum_{\omega : X(\omega) = x} \P(\{\omega\})$, agrupar os
termos dá $\E(X) = \sum_{\omega \in \Omega} \P(\{\omega\})\,X(\omega)$.
Então
\[
\E(X + Y) = \sum_{\omega} \P(\{\omega\})\bigl(X(\omega) + Y(\omega)\bigr)
= \sum_{\omega} \P(\{\omega\})X(\omega)
+ \sum_{\omega} \P(\{\omega\})Y(\omega) = \E(X) + \E(Y). \qedhere
\]
\end{proof}

\begin{definition}[Variáveis aleatórias independentes]\label{def:g12:sums:indep}
$X$ e $Y$ são \emph{independentes}\index{variável aleatória!independente}
se, para todos os valores $x, y$:
\[
\P(X = x \text{ e } Y = y) = \P(X = x)\,\P(Y = y).
\]
Várias variáveis $X_1, \dots, X_n$ são independentes se essa regra do
produto vale para toda escolha de valores de toda subfamília.
\end{definition}

\begin{proposition}\label{prop:g12:sums:prodexp}
Se $X$ e $Y$ são \omterm{def:g12:sums:indep}{independentes}, $\E(XY) = \E(X)\,\E(Y)$.
\end{proposition}

\begin{proof}
\[
\begin{aligned}
\E(XY) &= \sum_{x, y} xy\;\P(X = x \text{ e } Y = y)
= \sum_{x, y} xy\,\P(X=x)\P(Y=y) \\
&= \Bigl(\sum_x x\P(X=x)\Bigr)\Bigl(\sum_y y\P(Y=y)\Bigr). \qedhere
\end{aligned}
\]
\end{proof}

\begin{theorem}[Variância de uma soma]\label{thm:g12:sums:variance}
\index{variância!de uma soma}
Se $X$ e $Y$ são \emph{\omterm{def:g12:sums:indep}{independentes}}, então
\[
\V(X + Y) = \V(X) + \V(Y).
\]
Mais geralmente, para $X_1, \dots, X_n$ \omterm{def:g12:sums:indep}{independentes}:
$\V(X_1 + \dots + X_n) = \V(X_1) + \dots + \V(X_n)$.
\end{theorem}

\begin{proof}
Usando König--Huygens (\cref{prop:g12:randvar:konig}) e a linearidade:
\begin{align*}
\V(X + Y) &= \E\bigl((X+Y)^2\bigr) - \bigl(\E X + \E Y\bigr)^2\\
&= \E(X^2) + 2\E(XY) + \E(Y^2)
- \E(X)^2 - 2\E(X)\E(Y) - \E(Y)^2\\
&= \V(X) + \V(Y) + 2\bigl(\E(XY) - \E(X)\E(Y)\bigr),
\end{align*}
e o último colchete se anula para variáveis \omterm{def:g12:sums:indep}{independentes}
(\cref{prop:g12:sums:prodexp}). O caso geral segue por indução.
\end{proof}

\begin{example}[A binomial revisitada]\label{ex:g12:sums:binomial}
Uma variável binomial $X \sim \mathcal B(n, p)$ é uma soma
$X = X_1 + \dots + X_n$ de $n$ variáveis de Bernoulli \omterm{def:g12:sums:indep}{independentes}. Logo,
estruturalmente:
\[
\E(X) = np, \qquad \V(X) = np(1-p),
\]
o que recupera o \cref{thm:g12:randvar:binomial} sem cálculo nenhum.
\end{example}

\begin{proposition}[Média amostral]\label{prop:g12:sums:mean}
Sejam $X_1, \dots, X_n$ \omterm{def:g12:sums:indep}{independentes} e com a mesma \omterm{def:g12:randvar:rv}{distribuição} de $X$
(\omterm{def:g12:randvar:exp}{esperança} $\mu$, \omterm{def:g12:randvar:exp}{variância} $\sigma^2$), e seja
$M_n = \frac{X_1 + \dots + X_n}{n}$ sua \emph{média amostral}\index{média
amostral}. Então
\[
\E(M_n) = \mu, \qquad \V(M_n) = \frac{\sigma^2}{n}, \qquad
\sigma(M_n) = \frac{\sigma}{\sqrt n}.
\]
\end{proposition}

\begin{proof}
A linearidade dá $\E(M_n) = \frac{n\mu}{n} = \mu$; a \omterm{def:g12:condprob:indep}{independência} dá
$\V(X_1 + \dots + X_n) = n\sigma^2$, e dividir por $n$ multiplica a
\omterm{def:g12:randvar:exp}{variância} por $\frac{1}{n^2}$ (\cref{prop:g12:randvar:affine}).
\end{proof}

O $\frac{\sigma}{\sqrt n}$ é a \emph{lei da raiz quadrada}, fundamental:
para reduzir à metade as flutuações de uma \omterm{def:g11:stat:mean}{média}, quadruplique o tamanho
da \omterm{def:g10:proba:sample}{amostra}.

\begin{omfigure}
\begin{tikzpicture}
\begin{axis}[omaxis,
  width=9.6cm, height=5.6cm,
  xlabel={$n$}, ylabel={$\sigma(M_n)/\sigma$},
  xmin=0, xmax=105, ymin=0, ymax=1.15,
  xtick={4,16,64}, ytick={1,0.5,0.25},
  yticklabels={$1$,$\tfrac12$,$\tfrac14$}]
  \addplot[omDef, thick, domain=1:100, samples=80] {1/sqrt(x)};
  \draw[black, dashed, thin] (axis cs:0,0.5) -- (axis cs:4,0.5)
    -- (axis cs:4,0);
  \draw[black, dashed, thin] (axis cs:0,0.25) -- (axis cs:16,0.25)
    -- (axis cs:16,0);
  \draw[black, dashed, thin] (axis cs:64,0.125) -- (axis cs:64,0);
  \fill (axis cs:1,1) circle (1.5pt);
  \fill (axis cs:4,0.5) circle (1.5pt);
  \fill (axis cs:16,0.25) circle (1.5pt);
  \fill (axis cs:64,0.125) circle (1.5pt);
\end{axis}
\end{tikzpicture}

{\small A lei da raiz quadrada: cada quadruplicação do tamanho da \omterm{def:g10:proba:sample}{amostra}
apenas reduz à metade o \omterm{def:g11:stat:variance}{desvio padrão} da \omterm{def:g11:stat:mean}{média}.}
\end{omfigure}

\section{Desigualdades de concentração}

\begin{theorem}[Desigualdade de Markov]\label{thm:g12:sums:markov}
\index{desigualdade de Markov}
Se $X \geq 0$ e $a > 0$:
\[
\P(X \geq a) \leq \frac{\E(X)}{a}.
\]
\end{theorem}

\begin{proof}
Em $\E(X) = \sum_i p_i x_i$ (todos os termos não negativos), guarde apenas
os termos com $x_i \geq a$: cada um vale ao menos $a\,p_i$, logo
$\E(X) \geq a \sum_{x_i \geq a} p_i = a\,\P(X \geq a)$.
\end{proof}

\begin{theorem}[Desigualdade de Bienaymé--Chebyshev]\label{thm:g12:sums:chebyshev}
\index{desigualdade de Bienaymé--Chebyshev}
Para toda \omterm{def:g12:randvar:rv}{variável aleatória} $X$ e todo $\varepsilon > 0$:
\[
\P\bigl(\abs{X - \E(X)} \geq \varepsilon\bigr)
\leq \frac{\V(X)}{\varepsilon^2}.
\]
\end{theorem}

\begin{proof}
Aplique a desigualdade de Markov à variável não negativa
$Y = (X - \E(X))^2$ com $a = \varepsilon^2$:
\[
\P\bigl(\abs{X - \E(X)} \geq \varepsilon\bigr)
= \P(Y \geq \varepsilon^2)
\leq \frac{\E(Y)}{\varepsilon^2} = \frac{\V(X)}{\varepsilon^2}. \qedhere
\]
\end{proof}

\begin{omfigure}
\begin{tikzpicture}
\begin{axis}[omaxis,
  width=9.6cm, height=5cm,
  ylabel={}, xmin=-3.4, xmax=3.4, ymin=0, ymax=1.25,
  xtick={-1,0,1},
  xticklabels={$\E(X)-\varepsilon$,$\E(X)$,$\E(X)+\varepsilon$},
  ytick=\empty]
  \addplot[name path=bell, omDef, thick, domain=-3.2:3.2]
    {exp(-x^2/2)};
  \path[name path=xaxis] (-3.2,0) -- (3.2,0);
  \addplot[omDef!20] fill between[of=bell and xaxis,
    soft clip={domain=-1:1}];
  \addplot[omThm!20] fill between[of=bell and xaxis,
    soft clip={domain=-3.2:-1}];
  \addplot[omThm!20] fill between[of=bell and xaxis,
    soft clip={domain=1:3.2}];
  \draw[black, dashed, thin] (axis cs:-1,0) -- (axis cs:-1,0.75);
  \draw[black, dashed, thin] (axis cs:1,0) -- (axis cs:1,0.75);
\end{axis}
\end{tikzpicture}

{\small Concentração: Bienaymé--Chebyshev limita por
$\V(X)/\varepsilon^2$ a \omterm{def:g10:proba:distribution}{probabilidade} de $X$ cair nas caudas (em
vermelho), a distância pelo menos $\varepsilon$ de sua \omterm{def:g12:randvar:exp}{esperança}.}
\end{omfigure}

\begin{theorem}[Lei dos grandes números]\label{thm:g12:sums:lln}
\index{lei dos grandes números}
Sejam $X_1, \dots, X_n$ \omterm{def:g12:sums:indep}{independentes} e com a mesma \omterm{def:g12:randvar:rv}{distribuição}
(\omterm{def:g12:randvar:exp}{esperança} $\mu$, \omterm{def:g12:randvar:exp}{variância} $\sigma^2$) e $M_n$ sua \omterm{prop:g12:sums:mean}{média amostral}. Para
todo $\varepsilon > 0$:
\[
\P\bigl(\abs{M_n - \mu} \geq \varepsilon\bigr)
\leq \frac{\sigma^2}{n\,\varepsilon^2}
\xrightarrow[n \to +\infty]{} 0 .
\]
A \omterm{prop:g12:sums:mean}{média amostral} se concentra em torno da \omterm{def:g12:randvar:exp}{esperança}.
\end{theorem}

\begin{proof}
Bienaymé--Chebyshev aplicada a $M_n$, cuja \omterm{def:g12:randvar:exp}{esperança} é $\mu$ e cuja
\omterm{def:g12:randvar:exp}{variância} é $\frac{\sigma^2}{n}$ (\cref{prop:g12:sums:mean}).
\end{proof}

\begin{remark}
Este teorema é a ponte entre a teoria das \omterm{def:g10:proba:distribution}{probabilidades} e a estatística:
a \emph{\omterm{def:g10:stats:series}{frequência}} de um \omterm{def:g11:prob:model}{evento} ao longo de muitas repetições
\omterm{def:g12:sums:indep}{independentes} se aproxima de sua \emph{\omterm{def:g10:proba:distribution}{probabilidade}} (tome $X_i$ o
indicador do \omterm{def:g11:prob:model}{evento}, de modo que $\mu = p$). Ele justifica estimar uma
\omterm{def:g10:proba:distribution}{probabilidade} por simulação (\emph{método de Monte Carlo}) e uma proporção
populacional por uma pesquisa --- quantitativamente: veja o
\cref{ch:g12:contdist}.
\end{remark}

\begin{method}[Usar Bienaymé--Chebyshev]\label{met:g12:sums:chebyshev}
Para garantir $\P(\abs{M_n - \mu} \geq \varepsilon) \leq \alpha$, basta
que $n \geq \frac{\sigma^2}{\alpha\,\varepsilon^2}$. A cota é grosseira
(as flutuações reais costumam ser bem menores), mas perfeitamente geral:
nada exige da \omterm{def:g12:randvar:rv}{distribuição} além de uma \omterm{def:g12:randvar:exp}{variância} finita.
\end{method}

\section{Exercícios}

\begin{exercise}[$\star$]\label{exo:g12:sums:1}
Dois dados honestos são lançados; seja $S$ a soma. Usando a linearidade
(não a \omterm{def:g12:randvar:rv}{distribuição} de $S$!), calcule $\E(S)$; depois calcule $\V(S)$
usando a \omterm{def:g12:condprob:indep}{independência}, sabendo que um único dado honesto tem \omterm{def:g12:randvar:exp}{variância}
$\frac{35}{12}$.
\end{exercise}

\begin{exercise}[$\star$]\label{exo:g12:sums:2}
Seja $X \sim \mathcal B(100,\ 0.5)$. Limite $\P(X \geq 75)$ pela
desigualdade de Markov, e depois $\P(\abs{X - 50} \geq 25)$ por
Bienaymé--Chebyshev. Compare.
\end{exercise}

\begin{exercise}[$\star$]\label{exo:g12:sums:3}
Uma moeda honesta é lançada $n$ vezes e $F_n$ denota a \omterm{def:g10:stats:series}{frequência} de
caras. Quão grande deve ser $n$ para que, pela cota de
Bienaymé--Chebyshev,
$\P\bigl(\abs{F_n - 0.5} \geq 0.05\bigr) \leq 0.05$?
\end{exercise}

\begin{exercise}[$\star\star$]\label{exo:g12:sums:4}
Uma investidora reparte seu capital igualmente entre $n$ ativos
\omterm{def:g12:sums:indep}{independentes}, cada um com retorno esperado $\mu = 5\%$ e \omterm{def:g11:stat:variance}{desvio padrão}
$\sigma = 20\%$. Calcule a \omterm{def:g12:randvar:exp}{esperança} e o \omterm{def:g11:stat:variance}{desvio padrão} do retorno da
carteira $M_n$, e o número de ativos necessário para levar o \omterm{def:g11:stat:variance}{desvio padrão}
abaixo de $4\%$. Que princípio financeiro isso ilustra?
\end{exercise}

\begin{exercise}[$\star\star$]\label{exo:g12:sums:5}
Sejam $X$ e $Y$ \omterm{def:g12:sums:indep}{independentes}, ambas uniformes em $\{1, 2, 3\}$.
\begin{enumerate}
  \item Dê a \omterm{def:g12:randvar:rv}{distribuição} de $S = X + Y$ e calcule $\E(S)$, $\V(S)$
        diretamente a partir dela.
  \item Recupere os dois valores pela linearidade e pela aditividade da
        \omterm{def:g12:randvar:exp}{variância}.
\end{enumerate}
\end{exercise}

\begin{exercise}[$\star\star$]\label{exo:g12:sums:6}
Mostre que a aditividade da \omterm{def:g12:randvar:exp}{variância} pode falhar sem \omterm{def:g12:condprob:indep}{independência}:
calcule $\V(X + Y)$ para $Y = X$ e compare com $\V(X) + \V(Y)$. Para quais
variáveis $X$ vale $\V(2X) = 2\V(X)$?
\end{exercise}

\begin{exercise}[$\star\star$]\label{exo:g12:sums:7}
Suspeita-se que um dado esteja viciado. Ele é lançado $1200$ vezes e
mostra um seis $260$ vezes ($f = 0.2167$ em vez de
$\frac16 \approx 0.1667$). Sob a hipótese de que o dado é honesto, limite
$\P\bigl(\abs{F_n - \frac16} \geq 0.05\bigr)$ por Bienaymé--Chebyshev e
discuta se a hipótese de honestidade é plausível.
\end{exercise}

\begin{exercise}[$\star\star\star$]\label{exo:g12:sums:8}
\emph{(Uma desigualdade melhor para a moeda.)} Sejam
$X \sim \mathcal B(n,\ p)$ e $F_n = \frac Xn$.
\begin{enumerate}
  \item Mostre que $p(1-p) \leq \frac14$ para $p \in \intcc{0}{1}$.
  \item Deduza a cota válida para qualquer \omterm{def:g12:randvar:rv}{distribuição}
        $\P\bigl(\abs{F_n - p} \geq \varepsilon\bigr) \leq
        \frac{1}{4n\varepsilon^2}$.
  \item Quantas pessoas é preciso entrevistar para que a \omterm{def:g10:stats:series}{frequência}
        observada fique a menos de $3$ pontos da proporção verdadeira com
        \omterm{def:g10:proba:distribution}{probabilidade} de pelo menos $95\%$, usando essa cota? (As
        pesquisas reais usam estimativas mais finas, mas a ordem de
        grandeza está certa.)
\end{enumerate}
\end{exercise}

\section{Problema: a casa sempre ganha}

\begin{problem}[{Problema de fim de semana --- a lei dos grandes
números explica cassinos, pesquisas e seguros, e demole a falácia
do apostador no caminho}]
\label{pb:g12:sums:1}
Um jogador de roleta, depois de $100$ rodadas, está --- quase
tantas vezes quantas não --- \emph{ganhando}. O cassino que roda
um milhão de rodadas está ganhando com uma certeza que nenhum
tribunal questionaria. Mesmo jogo, mesma vantagem minúscula ---
a diferença é $\sqrt n$, e é o assunto deste capítulo
(\cref{prop:g12:sums:mean}, \cref{thm:g12:sums:chebyshev},
\cref{thm:g12:sums:lln}). Este problema faz a contabilidade do
cassino, dimensiona uma pesquisa eleitoral, precifica a
diversificação --- e desmonta a falácia mais cara da história do
jogo.

\medskip
\noindent\textbf{Parte I --- Fluência.}
\begin{enumerate}
  \item Dois dados \omterm{def:g12:sums:indep}{independentes}: calcule $V(X + Y)$
        (\cref{thm:g12:sums:variance}).
  \item Lance $100$ dados \omterm{def:g12:sums:indep}{independentes} e faça a \omterm{def:g11:stat:mean}{média} dos
        resultados: dê a \omterm{def:g12:randvar:exp}{esperança} e o \omterm{def:g11:stat:variance}{desvio padrão} da \omterm{prop:g12:sums:mean}{média
        amostral} (\cref{prop:g12:sums:mean}; para um dado,
        $\sigma = \sqrt{35/12} \approx 1.71$).
  \item Uma \omterm{def:g12:randvar:rv}{variável aleatória} não negativa tem \omterm{def:g12:randvar:exp}{esperança} $2$.
        O que a desigualdade de Markov
        (\cref{thm:g12:sums:markov}) diz sobre
        $\P(X \geq 10)$?
  \item Limite $\P\left(\abs{\bar X_{100} - 3.5} \geq
        0.5\right)$ para a \omterm{def:g11:stat:mean}{média} dos dados da questão 2 com
        Bienaymé--Chebyshev.
  \item As \omterm{def:g12:randvar:exp}{esperanças} sempre se somam; as \omterm{def:g12:randvar:exp}{variâncias}, só sob
        \omterm{def:g12:condprob:indep}{independência}: exiba $X, Y$ dependentes (dica:
        $Y = -X$) para os quais
        $V(X + Y) \neq V(X) + V(Y)$.
\end{enumerate}

\medskip
\noindent\textbf{Parte II --- A contabilidade da casa.}
Roleta europeia, aposta de $1$ euro no vermelho: ganha $+1$ com
\omterm{def:g10:proba:distribution}{probabilidade} $\frac{18}{37}$, perde $-1$ com \omterm{def:g10:proba:distribution}{probabilidade}
$\frac{19}{37}$.
\begin{enumerate}[resume]
  \item Calcule a \omterm{def:g12:randvar:exp}{esperança} e o \omterm{def:g11:stat:variance}{desvio padrão} do ganho de uma
        aposta.
  \item Um apostador faz $100$ apostas \omterm{def:g12:sums:indep}{independentes}; seja $G$
        o ganho total. Calcule $\E(G)$ e $\sigma(G)$, e depois a
        razão $\frac{\abs{\E(G)}}{\sigma(G)}$. O que uma deriva
        de um quarto de \omterm{def:g11:stat:variance}{desvio padrão} significa para as chances
        de o apostador estar ganhando ao fim da noite?
  \item O cassino vê $1\,000\,000$ de apostas. Calcule a
        \omterm{def:g12:randvar:exp}{esperança} e o \omterm{def:g11:stat:variance}{desvio padrão} de seu faturamento total e a
        mesma razão. Interprete o número $27$.
  \item Certifique com Chebyshev: limite a \omterm{def:g10:proba:distribution}{probabilidade} de o
        cassino \emph{perder} dinheiro ao longo do milhão de
        apostas.
  \item Sistemas de aposta: dobrar após as perdas, parar quando
        estiver ganhando\,\dots\ Usando a linearidade da
        \omterm{def:g12:randvar:exp}{esperança} (\cref{thm:g12:sums:linearity}) sobre a
        \omterm{def:g12:seq:sequence}{sequência} (possivelmente aleatória) de apostas, explique
        por que \emph{toda} estratégia em um jogo de vantagem
        negativa tem ganho esperado negativo --- o que um
        sistema vencedor violaria?
  \item Enuncie o que a lei dos grandes números
        (\cref{thm:g12:sums:lln}) promete sobre a \omterm{def:g10:stats:series}{frequência} do
        vermelho --- e o que ela \emph{não} promete sobre a
        próxima rodada depois de dez vermelhos seguidos. Nomeie
        a falácia.
  \item O ponto mais sutil: a lei funciona por \emph{diluição},
        não por compensação. Se os vermelhos ficam $100$ acima
        da expectativa depois de certa noite, esse excedente
        nunca é ``devolvido'' --- calcule o que acontece, em vez
        disso, com a \emph{fração} $\frac{100}{n}$ até
        $n = 10^6$, e reescreva o erro da falácia do apostador
        em uma frase.
\end{enumerate}

\medskip
\noindent\textbf{Parte III --- Pesquisas de opinião.}
\begin{enumerate}[resume]
  \item A partir do \cref{exo:g12:sums:8}: quantas pessoas é
        preciso entrevistar para que a \omterm{def:g10:stats:series}{frequência} observada
        fique a menos de $3$ pontos da verdade com \omterm{def:g10:proba:distribution}{probabilidade}
        de pelo menos $95\,\%$, pela cota de Chebyshev válida
        para qualquer \omterm{def:g12:randvar:rv}{distribuição}?
  \item Os institutos de pesquisa reais usam cerca de
        $1\,100$ pessoas para uma margem de $\pm 3$ pontos a
        $95\,\%$: a fórmula deles é
        $n \approx \frac{1.96^2 \times p(1-p)}
        {\varepsilon^2}$ com $p(1-p) \leq \frac14$. Avalie-a e
        explique a diferença em relação à questão 13 (o que
        Chebyshev não sabe sobre a forma das flutuações?).
  \item Para reduzir à metade a margem de erro, como a \omterm{def:g10:proba:sample}{amostra}
        deve crescer? De $\pm 3$ pontos com $1\,100$, que
        \omterm{def:g10:proba:sample}{amostra} dá $\pm 1.5$ ponto?
  \item O clássico contraintuitivo: o tamanho de \omterm{def:g10:proba:sample}{amostra}
        necessário nunca usou o tamanho da população ---
        $1\,100$ pessoas bastam para uma cidade ou para um
        continente. Aponte o lugar do modelo em que o tamanho
        da população está ausente e dê a analogia de cozinha
        que os institutos de pesquisa usam.
\end{enumerate}

\medskip
\noindent\textbf{Parte IV --- Diversificação.}
\begin{enumerate}[resume]
  \item A seguradora do \cref{pb:g12:randvar:1} ganha
        $30n$ euros em \omterm{def:g12:randvar:exp}{esperança} sobre $n$ apólices, com um
        \omterm{def:g11:stat:variance}{desvio padrão} dos sinistros de cerca de $315\sqrt n$.
        Para qual $n$ o lucro esperado finalmente supera um
        \omterm{def:g11:stat:variance}{desvio padrão} dos sinistros? O que o $\sqrt n$ está
        fazendo pela seguradora?
  \item O retorno anual de uma ação flutua com
        $\sigma = 20\,\%$. Reparta o dinheiro igualmente entre
        $25$ ações desse tipo \emph{\omterm{def:g12:sums:indep}{independentes}}: calcule o
        $\sigma$ da carteira. As finanças chamam a
        diversificação de o único almoço grátis --- em que
        consiste, exatamente, esse almoço?
  \item Agora suponha as $25$ ações perfeitamente
        correlacionadas (todas se movem juntas): qual é o
        $\sigma$ da carteira? Compare os dois extremos e diga
        que hipótese todo raciocínio de diversificação aluga em
        segredo --- e o que aconteceu quando ela falhou em todo
        o sistema em 2008.
  \item Finale --- a sinfonia do $\sqrt n$: as somas flutuam
        como $\sqrt n$ enquanto suas \omterm{def:g11:stat:mean}{médias} se assentam como
        $\frac{1}{\sqrt n}$; toque o tema pelas quatro
        indústrias deste problema (cassino, pesquisas, seguros,
        carteiras), nomeie os dois abusadores de plantão (a
        falácia do apostador e a \omterm{def:g12:condprob:indep}{independência} falsa) e dê o
        ponteiro adiante: a \emph{forma} das flutuações --- o
        sino --- é a estrela \omterm{def:g12:limcont:continuity}{contínua} do próximo capítulo, e
        seu teorema completo coroa os volumes de graduação.
\end{enumerate}
\end{problem}
